% Drawing 147/05-2 -- folder 147, batch 05, archive page 96 (PDF page 97).
% Transcription: le même dessin réduit à sa partie supérieure : un sommet portant deux boucles aplaties numérotées 2 et 3, avec le numéro 1 au-dessus du sommet ; deux arêtes parallèles vers un second sommet, puis une arête vers un sommet portant une boucle. La phrase s'interrompt après le dessin.
% Kind: graph with flattened loops (upper part of 05-1). Method 2 (schematic redraw from the facsimile). Status: ok.
% Notes: The upper part of the graph of 05-1: a vertex with the number 1 above it and two flattened loops numbered 2 and 3, two parallel edges to a second vertex, then an edge to a vertex carrying a loop. The loop is at the right edge of the sheet's text block; its size is approximate.
% Coordinates: 1 unit = 1 mm = 1/100 of the crop width (bbox in page fractions: [0.34, 0.821, 0.55, 0.869]).
% Redrawn by Claude (claude-opus-5-5), issue #73, 2026-11-07; not checked by a human; not compiled (no TeX in the session).
% Rights: reproduces part of the work -- behind RENDER_TRANSCRIPTS, never published by itself.
\documentclass[tikz,border=3pt]{standalone}
\usepackage{amsmath,amssymb}
\usetikzlibrary{patterns,arrows.meta}
\tikzset{pencil/.style={black!45}, struck/.style={black!35}, mark/.style={circle,draw,inner sep=1.3pt}, pt/.style={circle,fill,inner sep=0.9pt}, opt/.style={circle,draw,fill=white,inner sep=1.1pt}}
\begin{document}
\begin{tikzpicture}[x=1mm,y=1mm,line cap=round,line join=round,font=\scriptsize]
\node[pt] at (26.7,-16.2) {};
\draw (26.7,-16.2) -- (7.6,-10.5);
\draw (26.7,-16.2) -- (10.5,-25.7);
\node[inner sep=1pt] at (26.7,-8.6) {\tiny $1$};
\node[inner sep=1pt] at (5.7,-5.7) {\tiny $2$};
\node[inner sep=1pt] at (6.7,-22.9) {\tiny $3$};
\draw (26.7,-16.2) -- (49.5,-16.2);
\draw plot[smooth, tension=0.55] coordinates {(26.7,-16.2) (33.3,-21.9) (42.9,-21.9) (49.5,-16.2)};
\node[pt] at (49.5,-16.2) {};
\draw (49.5,-16.2) -- (73.3,-17.2);
\node[pt] at (73.3,-17.2) {};
\draw (84.8,-17.2) ellipse[x radius=11.4, y radius=5.7];
\end{tikzpicture}
\end{document}
